% GENERATED FILE. Edit canonical lessons, then run npm run generate:books.
% canonical-source-hash: fnv1a64:186c8e6e675e5c52
% canonical-lessons: MR-C539-prakasak, MR-C539-atmacaritra, MR-C539-pradarsan, MR-C539-digdarsak, MR-C539-dad-dene

\chapter{Books, Shows and Applause}
\label{ch:mr-books-shows-and-applause}
\hlchaptermodality{539}

\begin{chapteropening}
\textbf{By the end of this chapter:} \emph{I can talk about publishers, exhibitions, directors and applause.}

Complete the last lesson of chapter 539: I can talk about publishers, exhibitions, directors and applause.
\end{chapteropening}

\section[\texorpdfstring{\mr{प्रकाशक} (prakāśak)}{प्रकाशक (prakāśak)}]{\mr{प्रकाशक} (prakāśak) — a publisher}

\label{lesson:MR-C539-prakasak}



\begin{quote}
Before the new one: say the Marathi for to keep the beat, then the Marathi for to sway.
\end{quote}

\subsection*{You'll want to know: \mr{प्रकाशक}}
\textbf{\mr{प्रकाशक}} — \emph{prakāśak} — \textquotedblleft{}a publisher\textquotedblright{}.

\begin{quote}
\mr{प्रकाशकाने} \mr{माझ्या} \mr{मावशीचं} \mr{पुस्तक} \mr{छापलं}.
\end{quote}

\begin{quote}
\emph{prakāśakāne mājhyā māvśīcaṁ pustak chāplaṁ.}
\end{quote}

\begin{quote}
The publisher printed my aunt's book.
\end{quote}

Masculine.

\begin{cousinweb}
\textbf{\mr{प्रकाशक}} (\emph{prakāśak}) is a tatsama, taken from Sanskrit \emph{prakāśaka}, \textquotedblleft{}bringing to light\textquotedblright{} (Monier-Williams).
\end{cousinweb}

\begin{grammarlens}[title={What you've built}]
One more word for describing what you lived and felt.
\end{grammarlens}

\subsection*{Your turn}
\begin{itemize}
\raggedright
  \item \emph{Say it:} \emph{prakāśak}
  \item \emph{Say it:} \emph{prakāśak}, once more
  \item \emph{Say it:} \emph{prakāśakāne mājhyā māvśīcaṁ pustak chāplaṁ.} — The publisher printed my aunt's book.
  \item \emph{Recall:} say the Marathi for to keep the beat, then the Marathi for to sway, then say \emph{prakāśak} again
\end{itemize}

\begin{tcolorbox}[breakable,colback=teal!4,colframe=teal!35!black,title={Before you move on}]
What does \mr{प्रकाशक} mean? (\textquotedblleft{}A publisher\textquotedblright{}.) Say it once more.
\end{tcolorbox}
\section[\texorpdfstring{\mr{आत्मचरित्र} (ātmacaritra)}{आत्मचरित्र (ātmacaritra)}]{\mr{आत्मचरित्र} (ātmacaritra) — an autobiography}

\label{lesson:MR-C539-atmacaritra}



\begin{quote}
Before the new one: say the Marathi for to sway, then the Marathi for a publisher.
\end{quote}

\subsection*{You'll want to know: \mr{आत्मचरित्र}}
\textbf{\mr{आत्मचरित्र}} — \emph{ātmacaritra} — \textquotedblleft{}an autobiography\textquotedblright{}.

\begin{quote}
\mr{आजोबांनी} \mr{त्यांचं} \mr{आत्मचरित्र} \mr{लिहिलं}.
\end{quote}

\begin{quote}
\emph{ājobānnī tyāñcaṁ ātmacaritra lihilaṁ.}
\end{quote}

\begin{quote}
Grandfather wrote his autobiography.
\end{quote}

Neuter.

\begin{cousinweb}
\textbf{\mr{आत्मचरित्र}} (\emph{ātmacaritra}) is a tatsama compound of Sanskrit \emph{ātman}, \textquotedblleft{}self\textquotedblright{}, and \emph{caritra}, \textquotedblleft{}conduct, a life story\textquotedblright{} (Monier-Williams).
\end{cousinweb}

\begin{grammarlens}[title={What you've built}]
One more word for describing what you lived and felt.
\end{grammarlens}

\subsection*{Your turn}
\begin{itemize}
\raggedright
  \item \emph{Say it:} \emph{ātmacaritra}
  \item \emph{Say it:} \emph{ātmacaritra}, once more
  \item \emph{Say it:} \emph{ājobānnī tyāñcaṁ ātmacaritra lihilaṁ.} — Grandfather wrote his autobiography.
  \item \emph{Recall:} say the Marathi for to sway, then the Marathi for a publisher, then say \emph{ātmacaritra} again
\end{itemize}

\begin{tcolorbox}[breakable,colback=teal!4,colframe=teal!35!black,title={Before you move on}]
What does \mr{आत्मचरित्र} mean? (\textquotedblleft{}An autobiography\textquotedblright{}.) Say it once more.
\end{tcolorbox}
\section[\texorpdfstring{\mr{प्रदर्शन} (pradarśan)}{प्रदर्शन (pradarśan)}]{\mr{प्रदर्शन} (pradarśan) — an exhibition; a show}

\label{lesson:MR-C539-pradarsan}



\begin{quote}
Before the new one: say the Marathi for a publisher, then the Marathi for an autobiography.
\end{quote}

\subsection*{You'll want to know: \mr{प्रदर्शन}}
\textbf{\mr{प्रदर्शन}} — \emph{pradarśan} — \textquotedblleft{}an exhibition; a show\textquotedblright{}.

\begin{quote}
\mr{शहरात} \mr{चित्रांचं} \mr{मोठं} \mr{प्रदर्शन} \mr{भरलं} \mr{होतं}.
\end{quote}

\begin{quote}
\emph{śahrāt citrāñcaṁ moṭhaṁ pradarśan bharlaṁ hotaṁ.}
\end{quote}

\begin{quote}
A big exhibition of paintings was on in the city.
\end{quote}

Neuter.

\mr{प्रदर्शन} \mr{भरणे}, an exhibition is held.

\begin{cousinweb}
\textbf{\mr{प्रदर्शन}} (\emph{pradarśan}) is a tatsama, taken from Sanskrit \emph{pradarśana}, \textquotedblleft{}appearance, display\textquotedblright{} (Monier-Williams).
\end{cousinweb}

\begin{grammarlens}[title={What you've built}]
One more word for describing what you lived and felt.
\end{grammarlens}

\subsection*{Your turn}
\begin{itemize}
\raggedright
  \item \emph{Say it:} \emph{pradarśan}
  \item \emph{Say it:} \emph{pradarśan}, once more
  \item \emph{Say it:} \emph{śahrāt citrāñcaṁ moṭhaṁ pradarśan bharlaṁ hotaṁ.} — A big exhibition of paintings was on in the city.
  \item \emph{Recall:} say the Marathi for a publisher, then the Marathi for an autobiography, then say \emph{pradarśan} again
\end{itemize}

\begin{tcolorbox}[breakable,colback=teal!4,colframe=teal!35!black,title={Before you move on}]
What does \mr{प्रदर्शन} mean? (\textquotedblleft{}An exhibition; a show\textquotedblright{}.) Say it once more.
\end{tcolorbox}
\section[\texorpdfstring{\mr{दिग्दर्शक} (digdarśak)}{दिग्दर्शक (digdarśak)}]{\mr{दिग्दर्शक} (digdarśak) — a film director; a stage director}

\label{lesson:MR-C539-digdarsak}



\begin{quote}
Before the new one: say the Marathi for an autobiography, then the Marathi for an exhibition.
\end{quote}

\subsection*{You'll want to know: \mr{दिग्दर्शक}}
\textbf{\mr{दिग्दर्शक}} — \emph{digdarśak} — \textquotedblleft{}a film director; a stage director\textquotedblright{}.

\begin{quote}
\mr{दिग्दर्शकाने} \mr{नवीन} \mr{चित्रपटाची} \mr{घोषणा} \mr{केली}.
\end{quote}

\begin{quote}
\emph{digdarśakāne navīn citrapṭācī ghoṣṇā kelī.}
\end{quote}

\begin{quote}
The director announced a new film.
\end{quote}

Masculine.

\begin{cousinweb}
\textbf{\mr{दिग्दर्शक}} (\emph{digdarśak}) is a tatsama coinage from Sanskrit \emph{diś}, \textquotedblleft{}direction\textquotedblright{}, and \emph{darśaka}, \textquotedblleft{}one who shows\textquotedblright{} (Monier-Williams).
\end{cousinweb}

\begin{grammarlens}[title={What you've built}]
One more word for describing what you lived and felt.
\end{grammarlens}

\subsection*{Your turn}
\begin{itemize}
\raggedright
  \item \emph{Say it:} \emph{digdarśak}
  \item \emph{Say it:} \emph{digdarśak}, once more
  \item \emph{Say it:} \emph{digdarśakāne navīn citrapṭācī ghoṣṇā kelī.} — The director announced a new film.
  \item \emph{Recall:} say the Marathi for an autobiography, then the Marathi for an exhibition, then say \emph{digdarśak} again
\end{itemize}

\begin{tcolorbox}[breakable,colback=teal!4,colframe=teal!35!black,title={Before you move on}]
What does \mr{दिग्दर्शक} mean? (\textquotedblleft{}A film director; a stage director\textquotedblright{}.) Say it once more.
\end{tcolorbox}
\section[\texorpdfstring{\mr{दाद} \mr{देणे} (dād deṇe)}{दाद देणे (dād deṇe)}]{\mr{दाद} \mr{देणे} (dād deṇe) — to applaud; to show appreciation}

\label{lesson:MR-C539-dad-dene}



\begin{quote}
Before the new one: say the Marathi for an exhibition, then the Marathi for a film director.
\end{quote}

\subsection*{You'll want to know: \mr{दाद} \mr{देणे}}
\textbf{\mr{दाद} \mr{देणे}} — \emph{dād deṇe} — \textquotedblleft{}to applaud; to show appreciation\textquotedblright{}.

\begin{quote}
\mr{गायिकेच्या} \mr{गाण्याला} \mr{लोकांनी} \mr{मनापासून} \mr{दाद} \mr{दिली}.
\end{quote}

\begin{quote}
\emph{gāyikecyā gāṇyālā lokānnī manāpāsūn dād dilī.}
\end{quote}

\begin{quote}
People warmly applauded the singer's song.
\end{quote}

Past, with -\mr{ने} (he …): \mr{त्याने} \mr{दाद} \mr{दिली}, \emph{tyāne dād dilī}.

What is praised takes -\mr{ला}: \mr{गाण्याला} \mr{दाद} \mr{देणे}.

\begin{cousinweb}
\textbf{\mr{दाद} \mr{देणे}} (\emph{dād deṇe}) is \mr{दाद}, a loan from Persian \emph{dād}, \textquotedblleft{}justice, redress\textquotedblright{}, with \mr{देणे}: to give a performer their due.
\end{cousinweb}

\begin{grammarlens}[title={What you've built}]
One more word for describing what you lived and felt.
\end{grammarlens}

\subsection*{Your turn}
\begin{itemize}
\raggedright
  \item \emph{Say it:} \emph{dād deṇe}
  \item \emph{Say it:} \emph{dād deṇe}, once more
  \item \emph{Say it:} \emph{gāyikecyā gāṇyālā lokānnī manāpāsūn dād dilī.} — People warmly applauded the singer's song.
  \item \emph{Recall:} say the Marathi for an exhibition, then the Marathi for a film director, then say \emph{dād deṇe} again
\end{itemize}

\begin{tcolorbox}[breakable,colback=teal!4,colframe=teal!35!black,title={Before you move on}]
What does \mr{दाद} \mr{देणे} mean? (\textquotedblleft{}To applaud; to show appreciation\textquotedblright{}.) Say it once more.
\end{tcolorbox}
